El presente artículo científico expone el diseño, implementación y validación de un sistema de gestión de tickets orientado a optimizar la trazabilidad y resolución de incidencias en entornos organizacionales. Se propone una arquitectura híbrida que combina un portal web desarrollado en React con un backend en Laravel y una aplicación móvil para técnicos. El proyecto se fundamenta en metodologías ágiles, específicamente Scrum, lo que permitió entregas iterativas y retroalimentación continua. La investigación se apoya en múltiples artículos académicos que comparan frameworks, metodologías y arquitecturas (Artículo 1; Artículo 2; Artículo 3). Los resultados evidencian mejoras significativas en la centralización de la información, reducción de tiempos de atención y capacidad de personalización por roles. Finalmente, se discute la pertinencia de los frameworks seleccionados y las implicaciones de adoptar prácticas ágiles en el desarrollo de software moderno.

Palabras clave: Gestión de tickets; Soporte técnico; Metodologías ágiles; Laravel; React; React Native; Aplicación web; Aplicación móvil.